% EGA V, English translation by P. Blass, 1987 (cote 24): archive cote 24, site page 6 (PDF page 7).
% typescript. Draft compared: none.
% AI-transcribed by Claude Opus 5.5 (claude-opus-5-5, adaptive thinking, effort high), Message Batches API, batch ega24b (msgbatch_01RGJJeHdR96U16ZWFuUjtDx), 2026-11-13; not yet checked by eye.
\scanpage{24/6}
\mspage{6}

(N.B. to tell the truth the determination of the relative divisors
of P/S could be done with the means available right now, using
results of Ch. III and could be added as an example to Par. 19..).

Let us now make the base change $S^1 = P^V \to S$ and let us
consider the diagonal section (or ``generic section'') of $P^V{}_{S1}$ =
$\underline{P}(E^V{}_{S1})$ over $S^1$ : we find a closed subscheme $\underset{S}{\underline{H}}$ of $P_{S1} = P \times_S P^V$
which is called the incidence scheme between $\underline{P}$ and $\underline{P}^V$ defined by
the image ideal of the canonical homomorphism
$O_{\underline{P}}(-1) \otimes_S O_{\underline{P}^V}(-1) \to O_{\underline{P}\,\underset{S}{\times}\,\underline{P}^V}$

from which one sees that it is a projective fibration over $P^V$, and
by symmetry it is also a projective fibration over P.  Let us note
that one recovers the ``special'' hyperplanes $H_\xi$ (for $\xi$ a section of
$P^V$ over S) by starting out from the ``universal hyperplane'' $\underline{H}$ and by
taking its inverse image for the base change $S \xrightarrow{\xi} P^V$.

The same remark holds for every point of $P^V$ with values in an
arbitrary S-prescheme S' which (considered as a section of $\underline{P}_S{}'$ over
S') allows us to define an $H_\xi$ \hand{$\subset$} $P_S{}'$ ; the latter is nothing else
but the inverse image of H by the base change $S' \xrightarrow{\xi} \underline{P}^V$.  In what
follows we assume a prescheme X of finite type over S and an S
morphism $f{:}X \;\to P$.  One of the main objectives of this paragraph
(section) is to study for every hyperplane $H_\xi$ of P its inverse
image $\underline{Y}_\xi = f^{-1}(H_\xi) = X \underset{P}{\times} H_\xi$ and especially to relate
% NOTE: a white box covers the end of this last typed line, after "to relate"; the rest of the line is hidden.
